\documentclass[a4paper, 12pt, twoside]{article}
\usepackage{ifxetex}
\ifxetex
  \usepackage{fontspec}
\else
  \usepackage[T1]{fontenc}
  \usepackage[utf8]{inputenc}
\fi
\usepackage{babel}
\usepackage{lmodern}
\usepackage{graphicx}
\usepackage{enumitem}
\usepackage{lscape}
\usepackage{subcaption}
\usepackage{imakeidx}
\usepackage{tocbibind}
\usepackage[hidelinks]{hyperref}

\usepackage{xcolor}
\usepackage{listings}
\lstset{%
	basicstyle=\footnotesize\ttfamily\linespread{1}\selectfont,%
	numbers=left,%
	backgroundcolor=\color{lightgray},%
	showstringspaces=false,%
	breaklines=true%
}
\usepackage[newfloat]{minted}
\SetupFloatingEnvironment{listing}{}
\usemintedstyle{emacs}
\setminted{linenos, breaklines, tabsize=4, bgcolor=lightgray}

\usepackage{csquotes}
\usepackage[normalem]{ulem}

\usepackage[margin=2.5cm]{geometry}
\usepackage{setspace}
\setlength\parindent{1cm}
\onehalfspacing

\begin{document}

\begin{titlepage}
    \vspace*{5cm}
    \begin{center}
        \Huge \textbf{Architecture}
    \end{center}
    \vfill
\end{titlepage}

\tableofcontents
\clearpage

\subsection*{md2tex}
\addcontentsline{toc}{subsection}{md2tex}

\texttt{md2tex} is built around a layered architecture cleanly separating text parsing, abstract representations, escaping mechanisms, and output rendering.

\subsection*{Core Modules}
\addcontentsline{toc}{subsection}{Core Modules}

\begin{itemize}
    \item \textbf{\texttt{cli.py} \& \texttt{tui/}}: User interface layers routing interactions via Click (CLI commands) or Textual (Terminal User Interface). Both interface layers utilize the \texttt{md2tex/services.py} core to execute Markdown conversions reliably.
    \item \textbf{\texttt{services.py}}: The orchestration layer. It receives requests, standardizes path interactions, constructs \texttt{RenderOptions}, calls the parser and renderer, and injects templates.
    \item \textbf{\texttt{parser.py}}: A block-level, line-by-line Markdown parser. It scans strings and safely extracts raw block entities (like \texttt{CodeBlock}, \texttt{BlockQuote}, \texttt{Table}, \texttt{Heading}) into AST \texttt{IR Nodes}. 
    \item \textbf{\texttt{nodes.py}}: Simple structure layer outlining the abstract tree (\texttt{Table}, \texttt{Heading}, \texttt{CodeBlock}, \texttt{FootnoteDef}, \texttt{ListBlock}, etc.).
    \item \textbf{\texttt{escaping.py}}: Security translation module correctly escaping character artifacts (\texttt{\$}, \texttt{\%}, \texttt{\{}) \textit{after} blocks are parsed, but \textit{before} formatting runs.
    \item \textbf{\texttt{renderer.py}}: Main TeX compilation module processing \texttt{IR Nodes}. It handles each node type (e.g., \texttt{\_render\_table}) and recursively processes Inline Formatting via Regex (Italics, Hyperlinks, Strikethrough, etc.).
\end{itemize}

\subsection*{Flow of Data}
\addcontentsline{toc}{subsection}{Flow of Data}

\texttt{Markdown String -> Block Parser -> IR Ast (nodes.py) -> Renderer (escape -> format -> inject) -> TeX String}


\end{document}
